\documentclass[../../main.tex]{subfiles}
\begin{document}

\begin{center}
  {\fontsize{48pt}{58pt}\selectfont \textbf{1967}\normalsize\textbf{年度}}
\end{center}

\begin{table}[htb]
  \centering
  \caption{出題テーマと難易度}
  \begin{tabular}{|c|p{9cm}|c|}
    \hline
    \textbf{問題番号} & \textbf{テーマ} & \textbf{難易度} \\
    \hline
    第1問 & 三次方程式の実根の個数と範囲 & \\
    \hline
    第2問 & 複素数の二次方程式 & \\
    \hline
    第3問 & 双曲線の接線が作る三角形の面積 & \\
    \hline
    第4問 & ペル方程式の整数解と漸化式 & \\
    \hline
    第5問 & 新課程：位置ベクトルの一次結合／旧課程：三角形を二等分する線分の最短長 & \\
    \hline
    第6問 & 新課程：曲線の比較と囲まれる面積／旧課程：回転体 & \\
    \hline
  \end{tabular}
\end{table}

\begin{table}[htb]
  \centering
  \caption{解答}
  \renewcommand{\arraystretch}{1.4}
  \begin{tabular}{|c|c|p{9cm}|}
    \hline
    \textbf{問題番号} & \textbf{小問} & \textbf{答え} \\
    \hline
    第1問 & (1) & 3実根をもつ（証明問題） \\
    \cline{2-3}
    & (2) & 正の根はただ1つで$1<x<1+\dfrac{2}{k}$ \\
    \hline
    第2問 & (i) & $(x,y)=\left(\pm\dfrac{1}{\sqrt{2}},\pm\dfrac{1}{\sqrt{2}}\right)$（複号同順） \\
    \cline{2-3}
    & (ii) & $(x,y)=(-1,3),(1,1)$ \\
    \hline
    第3問 & (1) & $\dfrac{ab}{2}\sqrt{1-\dfrac{a^2}{x_1^2}}$ \\
    \cline{2-3}
    & (2) & $\dfrac{ab}{2}$ \\
    \hline
    第4問 & -- & 非負の整数解は$(2+\sqrt{3})^n=x_n+y_n\sqrt{3}$（$n=0,1,2,\ldots$）で与えられる \\
    \hline
    第5問 & 新課程(1) & $P$が$AB$上にあることと位置ベクトルが係数和$1$の一次結合になることは同値 \\
    \cline{2-3}
    & 新課程(2) & $p=q=r=0$のみ \\
    \cline{2-3}
    & 新課程(3) & $(l,m,n)=(1-x-y,x,y)$ \\
    \cline{2-3}
    & 旧課程 &  \\
    \hline
    第6問 & 新課程(1) & 交点は$(1,1)$のみ（証明問題） \\
    \cline{2-3}
    & 新課程(2) & $1+\dfrac{4}{\pi}+\dfrac{\pi}{8}\log 2$ \\
    \cline{2-3}
    & 旧課程 &  \\
    \hline
  \end{tabular}
\end{table}

\paragraph{試験形態}
6問．試験時間は要確認．

\end{document}
